\documentclass[dvipsnames]{article}

\usepackage[utf8]{inputenc}
\usepackage{fontspec}
\usepackage{polyglossia}
\usepackage{unicode-math}
\usepackage{csquotes}
\setdefaultlanguage{russian}
\setotherlanguage{english}

\setmainfont[SmallCapsFont=Latin Modern Roman Caps]{CMU Serif}
\setsansfont[SmallCapsFont=Latin Modern Roman Caps]{CMU Sans Serif}
\setmonofont[SmallCapsFont=Latin Modern Roman Caps]{CMU Typewriter Text}
% \setmathfont{Latin Modern Math}

\usepackage[letterpaper,top=2cm,bottom=2cm,left=3cm,right=3cm,marginparwidth=1.75cm]{geometry}

\usepackage{amsmath}
\usepackage{graphicx}
\usepackage{amsthm}
\usepackage{array}
\usepackage{fancyhdr}
\usepackage{bussproofs}
\usepackage[colorlinks=true, allcolors=blue]{hyperref}
\usepackage[backend=biber]{biblatex}
\usepackage{xcolor}

\usepackage{minted2}
\usemintedstyle{xcode}
\setminted{escapeinside=||}
\newminted[stella]{cpp}{escapeinside=||,frame=single,texcomments,fontsize=\small}
\newmintinline[haskell]{haskell}{fontsize=\scriptsize}
\newmintinline[java]{java}{fontsize=\scriptsize}
\newmintinline[python]{python}{fontsize=\scriptsize}
\newmintinline[code]{cpp}{fontsize=\small}
\newmintinline[smallcode]{cpp}{fontsize=\small}
\let\smallcodeOrig\smallcode
\renewcommand{\smallcode}[2][]{%
  \ifmmode
    \if\relax\detokenize{#1}\relax
      \mathchoice{\mbox{\normalfont\small\ttfamily #2}}{\mbox{\normalfont\small\ttfamily #2}}{\mbox{\normalfont\small\ttfamily #2}}{\mbox{\normalfont\small\ttfamily #2}}%
    \else
      \mathchoice{\mbox{\normalfont\small\ttfamily #2}}{\mbox{\normalfont\small\ttfamily #2}}{\mbox{\normalfont\small\ttfamily #2}}{\mbox{\normalfont\small\ttfamily #2}}%
    \fi
  \else
    \smallcodeOrig[#1]{#2}%
  \fi}

\usepackage{sourcecodepro}
\usepackage{mdframed}

\newcommand{\Fsub}{F$_{\text{\textbf{<:}}}$}
\newcommand{\vto}{%
    \mathrel{%
        \vdash\hspace*{-5pt}%
        \raisebox{2.65pt}{\scalebox{.33}{\(\blacktriangleright\)}}%
    }%
}

\addbibresource{../references.bib}

\begin{document}

\thispagestyle{fancy}
\lhead{\includegraphics[height=1cm]{../images/itmo_logo}}
\rhead{\includegraphics[height=1cm]{../images/stella-logo}}
\cfoot{Языки программирования. Семантика и системы типов (весна 2026, Университет ИТМО)}

\section*{Тема 13. Теоретическое задание \textcolor{blue}{(выполнил Павел Григоренко)}}%

\begin{enumerate}
  \item Предложите кодировку для типов-сумм~\cite[\S 11.9]{TaPL-ru} в Системе \Fsub~(аналогично кодировке пар~\cite[\S 26.3]{TaPL-ru}):
  \begin{enumerate}
    \item определите тип \smallcode{Sum T₁ T₂};
    \item выпишите типы и определения функций \smallcode{inl}, \smallcode{inr} и \smallcode{match},
    соответствующих конструкциям \smallcode{inl(t)}, \smallcode{inr(t)} и
    \smallcode{case t of inl(x₁)⇒t₁ | inr(x₂)⇒t₂};
    \item покажите (построив соответствующее дерево вывода), что выполняется правило
    \begin{prooftree}
      \AxiomC{$\Gamma \vdash \smallcode{S₁ <: T₁}$}
      \AxiomC{$\Gamma \vdash \smallcode{S₂ <: T₂}$}
      \RightLabel{\textsc{S-Sum}}
      \BinaryInfC{$\Gamma \vdash \smallcode{Sum S₁ S₂  <:  Sum T₁ T₂}$}
    \end{prooftree}
  \end{enumerate}

  \color{blue}
  \textbf{Решение.}
  \emph{место для решения}
  \color{black}

  \clearpage

  \item Для каждого из следующих типов ядерной Системы \Fsub~\cite[Рис.~26-1]{TaPL-ru}
  постройте (любой) замкнутый терм этого типа вместе с деревом вывода типа
  (включая поддеревья вывода подтипизации, где необходимо)
  или покажите, что это невозможно:
  \begin{enumerate}
    \item \smallcode{∀X<:Top. Top → X}
    \item \smallcode{∀X<:Top. ∀Y<:X. X → Y}
    \item \smallcode{∀X<:Top. ∀Y<:X. ∀Z<:Y. Z → X}
  \end{enumerate}

  \color{blue}
  \textbf{Решение.}
  \emph{место для решения}
  \color{black}

  \clearpage

  \item Постройте дерево вывода, соответствующее алгоритму проверки типов для ядра \Fsub~\cite[Рис.~28-3]{TaPL-ru},
  для следующих замкнутых термов:
  \begin{enumerate}
    \item \smallcode{λX<:Nat→Nat. λy:X. y 0  :  ∀X<:Nat→Nat. X → Nat}
    \item \smallcode{(λX<:{a:Nat}. λx:X. succ x.a) [{a:Nat,b:Bool}] {a=0,b=true}  :  Nat}
    \item \smallcode{(λX<:Top. λx:X. x) [Nat→Nat] (λn:Nat. succ n)  :  Nat→Nat}
  \end{enumerate}

  \color{blue}
  \textbf{Решение.}
  \emph{место для решения}
  \color{black}

  \clearpage

  \item Для контекста
  $\Gamma = \smallcode{X<:Top, Y<:Nat→Bool, Z<:Y, W<:Z, B<:X→Nat, C<:B}$
  укажите, какие из следующих утверждений о выявлении верны.
  Для неверных поясните, какой наименьший надтип получается вместо указанного:
  \begin{enumerate}
    \item $\Gamma \vto \smallcode{Top} \Uparrow \smallcode{Top}$
    \item $\Gamma \vto \smallcode{Y} \Uparrow \smallcode{Top}$
    \item $\Gamma \vto \smallcode{Z} \Uparrow \smallcode{Nat→Bool}$
    \item $\Gamma \vto \smallcode{W} \Uparrow \smallcode{Top}$
    \item $\Gamma \vto \smallcode{B} \Uparrow \smallcode{X→Nat}$
    \item $\Gamma \vto \smallcode{C} \Uparrow \smallcode{Top}$
  \end{enumerate}

  \color{blue}
  \textbf{Решение.}
  \emph{место для решения}
  \color{black}
\end{enumerate}

\section*{Список литературы}
\printbibliography[heading=none]

\end{document}
